\documentclass[10pt]{article}

\usepackage[utf8]{inputenc}

\usepackage[T1]{fontenc}

\usepackage{amsmath}

\usepackage{amsfonts}

\usepackage{amssymb}

\usepackage[version=4]{mhchem}

\usepackage{stmaryrd}

\usepackage{graphicx}

\usepackage[export]{adjustbox}

\graphicspath{ {./images/} }

\usepackage{bbold}



\begin{document}

УСТОЙЧИВОЕ РАСПРЕДЕЛЕНИЕ - распределение вероятностей, обладающих свойством, что для любых $a_{1}>0, b_{1} \cdot a_{2}>0, b_{2}$ пимеет место соотногіеніе



\[

F\left(a_{1} x+b_{1}\right) * F\left(a_{2} x+b_{2}\right)=F(a x+b)

\]



где $a>0$ и $b$ - нек-рые постоянные, $F$ - функция распределения У. р., *- символ операции свертки двух функциӥ распределения.



Характеристическая функция У. р.



\[

\varphi(t)=\exp \left\{i d t+c|t|^{\alpha}\left[1+i \beta \frac{t}{|t|} \omega(t, \alpha)\right]\right\} \text {, }

\]



где $0<\alpha \leqslant 2,-1 \leqslant \beta \leqslant 1, c \geqslant 0, d$ - любое действительное число, и



\[

\omega(t, \alpha)=\left\{\begin{array}{lll}

\operatorname{tg} \frac{\pi \alpha}{2}, & \text { при } & \alpha \neq 1, \\

\frac{2}{\pi} \ln |t|, & \text { при } & \alpha=1 .

\end{array}\right.

\]



Число $\alpha$ наз. по о а $а$ а т е ле м у с т о й ч и в ого р а с п р е д е лен и я. У. р. с показателем $\alpha=2$ является норлальное распределение, примером У. р. с показателем $\alpha=1$ служит Коши распределение, У. р. является вырожденное распределение на прямой, У. р.безгранично делимое распределение, и У. р. с показателем $\alpha, 0<\alpha<2$, имеет Леви каноническое представление с характеристиками $\sigma^{2}=0$,



\[

\begin{gathered}

M(x)=c_{1} /(x)^{\alpha}, \quad N(x)=-c_{2} / x^{2}, \\

c_{1} \geqslant 0, \quad c_{2} \geqslant 0, \quad c_{1}+c_{2}>0,

\end{gathered}

\]



$\vartheta$ - любое действительное число.



Для У. р., за исключением вырожденного распределения, существуют плотности. Эти плотности бесконечно дифференцируемы, одновершинны и отличны от нуля или на всей прямой, или на полупрямой. Для У. р. с показателем $\alpha, 0<\alpha<2$, при $\delta<\alpha$ выполняются coотношения



\[

\int_{-\infty}^{\infty}|x|^{\delta} p(x) d x<\infty, \quad \int_{-\infty}^{\infty}|x|^{\alpha} p(x) d x=\infty

\]



$p(x)$ - плотность У. p. Явный вид плотностеӥ $\mathrm{У} . \mathbf{p}$. известен лишь в немногих случаях. Одной из основных задач теории У. р. является описание их притяжения областей.



В совокупности У. р. выделяется класс с т р о г о у с т о й ч и в ы $\mathbf{p}$ а с п р е д е л е н и ӥ, для к-рых имеет место равенство (1) при $b_{1}=b_{2}=b=0$. Характеристич. функции строго У. р. с показателем $\alpha(\alpha \neq 1)$ даются формулой (2) при $d=0$. При $\alpha=1$ строго У. р. является лишь распределение Коши. Спектрально положительные (отрицательные) У. р. характеризуются тем, что в канонич. представлении Леви $M(x) \equiv 0(N(x) \equiv$ $\equiv 0$ ). Для спектрально положительных У. р. существует преобразование Лапласа при $\operatorname{Re} s \geqslant 0$ :



$\int_{-\infty}^{\infty} e^{-s x} p(x) d x=\left\{\begin{array}{lll}\exp \left\{-c s^{\alpha}-d s\right\}, & \text { при } & \alpha<1, \\ \exp \{c s \ln s-d s\}, & \text { при } & \alpha=1, \\ \exp \left\{c s^{\alpha}-d s\right\}, & \text { при } & \alpha>1,\end{array}\right.$



где $p(x)$-плотность спсктрально положительного У. p. с показателем $\alpha, 0<\alpha<2, c>0, d$-действительное число, у многозначных функций $\ln s ; s^{\alpha}$ выбираются те ветви, для к-рых $\ln s$ действительный, а $s^{\alpha}>0$ при $s>0$.



У. р., как безгранично делимому распределению, соответствует однородный случайный процесс с независимыми приращениями. Стохастически неирерывный однородный случайный процесс с независимыми приращениями $\{x(\tau), \tau \geqslant 0\}$ наз. у с т о й ч и в ы м, если приращение $x(1)-x(0)$ имеет $\mathrm{Y} . \mathrm{p}$.



Jum.: [1] Гнед е н ко Б. В., К олмогоров А. Н., Предельные распределения для сумм независимых случайных ве-



\includegraphics[max width=\textwidth, center]{2023_07_09_6b616bf1b681ed8a362dg-1(1)}

н о в Ю А., Теория вероятностей, 2 изд., М., 1973; [3] и 6 р аг и м о в И. А., Л и н н и к Ю. В., Независимые и стационарно связанные величины, М , 1965; [4] С к о р о х о д А. В., Случайные процессы с независимыми приращениями, М , 1964; [5] 3 о лот а Р е в В. М., Одномерные устойчивые растределения,



УСТОИЧИВОСТИ КРИТЕРИИ -- необходІмые І достаточные условия отрицательности действительных частей всех корней уравнения



\[

\lambda^{n}+a_{1} \lambda^{n-1}+\ldots-a_{n}=0 .

\]



У. к. используотся, когда применяется теорема ЈКпунова об устойчивости Іо первому приблиякению неподвижной точки автономной системы диффференциальных



\includegraphics[max width=\textwidth, center]{2023_07_09_6b616bf1b681ed8a362dg-1(2)}

употребителен следующий $У$. к., наз. к р и т е р и е м $\mathrm{P}$ а у с а $-\Gamma$ у р в ц а или к р и т е р и е м $\Gamma$ у рв и ц а: для отрицательности действительных частей всех корней уравнения (i) необходимо и достаточно, чтобы выполняласг, совокупность, неравенств $\Delta_{i}>0$, $i \in\{1, \ldots, n\}$, где



\[

\Delta_{1}=a_{1}, \quad \Delta_{2}=\left|\begin{array}{ll}

a_{1} & 1 \\

a_{3} & a_{2}

\end{array}\right|, \quad \Delta_{3}=\left|\begin{array}{lll}

a_{1} & 1 & 0 \\

a_{3} & a_{2} & a_{1} \\

a_{5} & a_{4} & a_{3}

\end{array}\right|, \ldots

\]



\begin{itemize}

  \item главные диагональные миноры матрицы

\end{itemize}



\begin{center}

\includegraphics[max width=\textwidth]{2023_07_09_6b616bf1b681ed8a362dg-1}

\end{center}



(на главной диагонали әтой матриңы стоят $a_{1}, \ldots, a_{n}$; при $i>n$ полагают $\left.a_{i}=0\right)$.



П ри $n=2$ У. к. Рауса - Гурвица выглядит особенно просто: для отрицательности действительных частей корней уравнения $\lambda^{2}+a_{1} \lambda+a_{2}=0$ необходимо и достаточно, чтобы коәффициенты уравнения были положітельны: $a_{1}>0, a_{2}>0$.



При веяком $n \in \mathbb{N}$ для отрицательности действительных частей всех корней уравнения (*) необходимо (но при $n>2$ недостаточно), чтобы все коэффициенты уравнения были положительны: $a_{i}>0, i \in\{1, \ldots, n\}$. Если хоть один из оп ределителей $\Delta_{i}, i \in\{1, \ldots, n\}$, отрицателен, то у уравнения $(*)$ найдется корень, действительная часть к-рого положительна (это утверждение используется при применении теоремы Ляпунова о неустойчивости по первому приближенио неподвижної точки автономной системы дифференциальных уравнений, см. Устойчивость по Ліпунову). Если $\Delta_{i} \geqslant 0$ пля любого $i \in\{1, \ldots, n\}$, но $\Delta_{i}=0$ для некоторого $i \in\{1, \ldots, n\}$, то расположение корней уравнения $(*)$ относительно мнимой осіг тоже можно выяснить, не находя самих корней (см. [5], [8] гл. XVI, § 8).



Более простым для применения является к ритер и й Ль е а р а - III и п а $\mathrm{p}$ а: для отрицательности действительных частей всех корней уравнения $(*)$ необходимо и достаточно, чтобы выполнялась совокупность неравенств $a_{i}>0, i \in\{1, \ldots, n\}, \Delta_{n-2 i+1}>0$, $i \in\left\{1, \ldots,\left[\frac{n}{2}\right]\right\}$ (определители $\Delta_{i}$-те же, что в крітерии Рауса- Гурвига).



К р и т е р и і $Э$ р м и т а (исторически первый, см. [1], [10] § 3.1) позволяет с помощью конечного числа арифметич. дейіствий над коэфффициентами уравнения (*) определить, все ли корни әтого уравнения имеют отрицательную действительную часть. Сформулированный выше критерий Рауса - Гурвица - это модификация критерия Әрмита, найденная А. Гурвицем. Известен также У. к. Ляпунова (см. [3], [8] гл. XVI, § 5, [10] $\S 3.5)$.



При исследовании уетойчивости неподвижной точки диффференцируемого отображения (автономной системы





\end{document}